\section{Conclusão}

\paragraph{}
As atividades realizadas permitiram aplicar técnicas fundamentais de
reconhecimento e enumeração utilizadas em testes de penetração. Por meio da
identificação de portas abertas, serviços em execução, versões de software e
recursos disponíveis na aplicação web, foi possível obter uma visão mais ampla
da superfície de ataque do ambiente analisado.

\paragraph{}
A utilização de ferramentas como Nmap, cURL, Dirb, Gobuster, Sublist3r e
Metasploit demonstrou a importância da fase de coleta de informações para o
planejamento das etapas subsequentes de um pentest. As informações obtidas
permitem identificar potenciais vetores de ataque, validar exposições
desnecessárias e direcionar análises mais aprofundadas.

\paragraph{}
Os resultados evidenciam que a enumeração é uma etapa essencial do processo de
segurança ofensiva, pois fornece subsídios para a identificação de
vulnerabilidades, avaliação de riscos e fortalecimento da postura de segurança
dos ambientes analisados.
